\section{Conclusion and Discussion}
\label{sec:discussion}

We introduced iPerturb, which builds a cell-line-specific regulatory network from five evidence tiers and propagates perturbations through it with a learnable transfer function that follows the Hill function, so that every prediction can be traced to named genes and edges. On held-out CRISPRi targets in K562 and RPE1 the tiered network carries predictive signal that matched rewired and data-derived graphs do not, multi-hop propagation of training responses over the tiers, integrated across tiers, is the strongest read-out (\datanum{g12_k562_prop_integrated__nnls} and \datanum{g12_rpe1_prop_integrated__nnls} in $\rho_c$ against \datanum{k562_b2_rc} and \datanum{rpe1_b2_rc} for the linear baseline), and the Hill-form operator reproduces exact Hill propagation to \datanum{eq_cheb16_spmm_dx} while its matrix-product form is memory-lean on dense graphs. Two limitations are central. First, the tiers that carry the predictive gain (curated regulation and protein association) are the same for every cell line; the cell-line-specific enhancer tier, testable only in K562, differs strongly between lines but did not improve prediction over a foreign line's tier or over the base graph, so its context specificity is descriptive rather than predictive in this setting. Second, on the curated signed core, which reaches only \datanum{k562_signed_out_pct}--\datanum{rpe1_signed_out_pct}\% of the perturbed genes in these essential-gene screens, a nonlinear Hill transfer does not predict better than a linear one, and the gain over a linear low-rank model comes from association evidence rather than from curated regulation. The mechanism-level test sharpens this: the signed network predicts the direction of direct regulatory effects, but the signed multi-hop cascade is not supported beyond the first hop in these endpoint data, so we present longer paths only as candidate mechanisms, and we prove only that the multi-hop quantity is well defined and has a unique convergence point (Appendix~\ref{app:convergence}), not that it is correct. We read the coverage limitation as a statement about the screens, not about biology: the screens perturb mostly non-TF genes, and a signed core cannot propagate what it does not reach.

Several limitations qualify these results. The predicted endpoint is the mean response, not a distribution of cells. RNA abundance is a proxy for regulator activity and endpoint data cannot identify time constants, so the fixed-point operator is a steady-state model; neighbour attributions are tested by ablation but are association paths, not validated mechanisms. The folds and seeds reuse the same biological screens, and graph and no-graph models were tuned on comparable but not identical budgets; matched controls share their grid. The largest gains are obtained for genes in well-characterised complexes (for example, ribosomal-protein neighbours), which are enriched in essential-gene screens, and the transfer to screens with few such complexes is untested. The GeneHancer tier is empty for RPE1, so the cell-line-specific tier is tested in K562 only; a test with the training-perturbation tier (Level 5) is left to future work. Future work will add more cell lines and a genome-wide screen with more transcription-factor targets, and test edge-specific learnable shapes, which the matrix-multiplication form now makes affordable, on data in which a signed core has more reach.

\section*{Conflict of interests, acknowledgments, funding, and data availability}
\textit{Conflict of Interests.} The authors declare that they have no competing interests.

\textit{Funding.} This work was supported by the AI Institute of Fulbright University Vietnam and Google.

\textit{Data and Software Availability.} The CRISPRi Perturb-seq datasets analysed in this study are publicly available from Replogle et al.~\cite{replogle2022mapping}. Prior biological knowledge was obtained from TRRUST, CollecTRI, OmniPath, GeneHancer v5.26, STRING v12 and COXPRESdb; GeneHancer is licence-gated and is not redistributed. The source code, the frozen pre-registration of the evaluation protocol with its later addenda, figure-generation scripts and the files from which every number in this paper is generated accompany the submission and will be released in the public iPerturb repository (\url{https://github.com/kagtgi/iPerturb}) at publication.
